\section{Estimation Algorithm and Response Diagnostic}
\label{sec:estimation}

Let $m=|\Omega|$ and let
$\theta=(\cX,\{\bA^{(n)}\}_{n=1}^{N},\bbeta)$ collect all trainable
parameters. For an active response dictionary of size $Q+1$, we minimize
\begin{equation}
\begin{aligned}
  \mathcal{F}_Q(\theta)
  =
  \frac{1}{2|\Omega|}
  \norm{\Omega \od \left(\cY - f_{\bbeta}(\cS)\right)}_F^2
  +
  \frac{\lambda_X}{2}\norm{\cX}_F^2 \\
  \quad
  +
  \sum_{n=1}^{N}\frac{\lambda_n}{2}\norm{\bA^{(n)}}_F^2
  +
  \frac{\lambda_\beta}{2}\norm{\bbeta}_2^2,
\end{aligned}
  \label{eq:estimation_objective}
\end{equation}
where $f_{\bbeta}(s)=\sum_{q=0}^{Q}\beta_q\psi_q(s)$ and
$\{\psi_q\}_{q=0}^{Q}$ is the selected response dictionary. The monomial
basis $\psi_q(s)=s^q$ gives \eqref{eq:polynomial_linked_tucker}; in
implementation, the same polynomial space may be represented by centered,
Chebyshev, or otherwise conditioned bases. Non-polynomial choices such as RBF
or spline dictionaries use the same objective as long as the response remains
linear in $\bbeta$. The experiments use the monomial power dictionary. The
estimation alternates between an exact small response-coefficient update and a
descent step for the Tucker core and factors.

\subsection{Response Refresh as a Ridge Subproblem}

Fix the current Tucker signal $\cS$ and list the observed entries as
$y_\Omega\in\R^m$ and $s_\Omega\in\R^m$. Define the feature matrix
$\bPhi_{\Omega,Q}(\cS)\in\R^{m\times(Q+1)}$ by
\begin{equation}
  [\bPhi_{\Omega,Q}(\cS)]_{j,q}=\psi_q([s_\Omega]_j),
  \qquad
  q=0,\ldots,Q.
\end{equation}
The response update is the ridge regression
\begin{equation}
  \bbeta^+
  =
  \argmin_{\bbeta\in\R^{Q+1}}
  \frac{1}{2m}
  \norm{y_\Omega-\bPhi_{\Omega,Q}(\cS)\bbeta}_2^2
  +
  \frac{\lambda_\beta}{2}\norm{\bbeta}_2^2,
  \label{eq:link_refresh_problem}
\end{equation}
with closed-form solution
\begin{equation}
  \bbeta^+
  =
  \left(
  \frac{1}{m}\bPhi_{\Omega,Q}^{T}\bPhi_{\Omega,Q}
  +\lambda_\beta \bI
  \right)^{-1}
  \frac{1}{m}\bPhi_{\Omega,Q}^{T}y_\Omega.
  \label{eq:link_refresh_solution}
\end{equation}
This step solves only a $(Q+1)\times(Q+1)$ system. It is therefore cheap
relative to reconstructing or differentiating through the Tucker tensor, and
because it solves the $\bbeta$-subproblem exactly, it never increases
$\mathcal{F}_Q$ for the current Tucker signal.

\subsection{Descent for the Tucker Parameters}

With $\bbeta$ fixed, the residual is backpropagated through the entrywise response
and then through the Tucker map. The gradient with respect to the Tucker signal
is
\begin{equation}
  \bG_{\cS}
  =
  \frac{1}{m}
  \Omega\od
  \left(f_{\bbeta}(\cS)-\cY\right)
  \od f_{\bbeta}'(\cS).
  \label{eq:latent_gradient}
\end{equation}
The Tucker gradients are standard reverse contractions:
\begin{align}
  \nabla_{\cX}\mathcal{F}_Q
  &=
  \bG_{\cS}
  \times_1 (\bA^{(1)})^T
  \cdots
  \times_N (\bA^{(N)})^T
  +\lambda_X\cX, \\
  \nabla_{\bA^{(n)}}\mathcal{F}_Q
  &=
  [\bG_{\cS}]_{(n)}
  \bZ_n^T
  +
  \lambda_n\bA^{(n)},
  \label{eq:factor_gradient}
\end{align}
where
$\bZ_n=(\bA^{(N)}\otimes\cdots\otimes\bA^{(n+1)}\otimes
\bA^{(n-1)}\otimes\cdots\otimes\bA^{(1)})\cX_{(n)}^T$.
The Tucker update can use a line-search gradient step, a block coordinate
gradient step, or an adaptive first-order implementation, provided the step
satisfies the sufficient decrease condition used in Section~\ref{sec:theory}.

\subsection{Dictionary Continuation and Normalization}

Large or high-order response dictionaries can amplify Tucker scale ambiguity. For
nested dictionaries, such as the power or Chebyshev bases, we use continuation:
training starts with the Tucker-compatible identity response and activates
additional channels gradually. When a new channel is introduced, its coefficient
is initialized at zero, so the prediction is unchanged at the moment of
expansion. Non-nested dictionaries such as RBF and spline bases are trained with
the full selected dictionary from the start. After Tucker updates, factor
columns are rescaled and the inverse scales are absorbed into the core. This
fixes a numerical gauge without changing the represented signal $\cS$ or the
objective value.

\subsection{Alternating Recovery Scheme}

One outer iteration consists of:
\begin{enumerate}
  \item fixing the current Tucker signal and refreshing $\bbeta$ by
  \eqref{eq:link_refresh_solution};
  \item updating $\cX$ and $\{\bA^{(n)}\}$ by a descent step using
  \eqref{eq:latent_gradient}--\eqref{eq:factor_gradient};
  \item normalizing the Tucker gauge and, when the dictionary is nested,
  increasing the active dictionary size according to a continuation schedule.
\end{enumerate}
This procedure treats the response as part of the fitted representation: the
Tucker update is computed through the current nonlinear measurement map.

\subsection{Response Diagnostic}

The diagnostic measures whether the residual of a linear Tucker fit is
coherent with an entrywise response. Let $\cS_T$ be a fitted Tucker reconstruction
and let $\bbeta_T=(0,1,0,\ldots,0)$ denote the identity response. The frozen-Tucker
diagnostic is
\begin{align}
  E_{\mathrm{lin}}
  &=
  \norm{y_\Omega-\bPhi_{\Omega,1}(\cS_T)\bbeta_T}_2^2, \\
  E_{\mathrm{link}}
  &=
  \min_{\bbeta}
  \left[
  \norm{y_\Omega-\bPhi_{\Omega,Q}(\cS_T)\bbeta}_2^2
  +m\lambda_\beta\norm{\bbeta}_2^2
  \right], \\
  \Dlink
  &=
  \frac{E_{\mathrm{lin}}-E_{\mathrm{link}}}{E_{\mathrm{lin}}}.
  \label{eq:dlink}
\end{align}
Large $\Dlink$ means that refitting only the response coefficients can already
remove substantial residual energy without changing the Tucker fit. We use this
score as a model-selection and interpretation tool: high diagnostic yield
suggests that the nonlinear response is aligned with the residual structure,
whereas low diagnostic yield suggests that increasing multilinear rank or
changing the tensor prior may be more appropriate.
